% ----------------------
\subsection{Section 119 (8 July 1838)}
% ----------------------
	\label{DC119}
	
	% -----
	\subsubsection{Historical Background}
	% -----
		\label{DC119-HB}
		
		% --
		\paragraph{JSP}
		% --
		
			``On Sunday, 8 July 1838, JS dictated five revelations, each of which concerned church leadership or finances; one of these revelations outlined a plan for raising church revenue. The subject of church finances was not new to the Saints. From the time the church was established, JS dictated revelations and instituted programs related to economic and social concerns. An 1831 revelation on `the Laws of the Church of Christ' directed the Latter-day Saints to consecrate their property to the church bishop and then manage stewardships of property or other responsibilities assigned to them. Church members attempted to follow this program of consecration and stewardship in Jackson County, Missouri, but their attempts ended when they were driven out of the county in 1833. In the early and mid-1830s, JS and other church leaders in Ohio engaged in a number of business and banking ventures, most of which ultimately failed. The troubled situation of church finances was compounded by the nationwide panic of 1837 and the ensuing economic recession. During this period, JS and Sidney Rigdon incurred several thousand dollars of debt.
			
			``In the latter half of 1837, the bishops in Missouri and Kirtland, Ohio, took new steps to address the church's financial problems. In September, the church published an appeal from Bishop Newel K. Whitney and his counselors in Kirtland, calling on church members everywhere to `bring their tithes into the store house' to relieve church debts and to help establish the community of Saints in Missouri. While this general request did not include recommended donation amounts, in December 1837 a committee composed of Edward Partridge, the bishop of Zion; Isaac Morley, the first counselor in the bishopric; and John Corrill, the appointed `keeper of the Lord's Storehouse,' proposed that every head of household be asked to annually donate a certain percentage of net worth, with the percentage based on church needs for the year. To cover anticipated church expenses for 1838, the committee proposed a `tithing' of 2 percent. The committee believed that such a program would `be in some degree fullfilling the law of consecration.' In February 1838, when Thomas B. Marsh, the pro tempore president of Zion, wrote to JS about coming to Missouri, Marsh reported that the Saints there `seem to wish to have the whole law of God lived up to; and we think that the church will rejoice to come up to the law of consecration, as soon as their leaders shall say the word, or show them how to do it.' In April a revelation called for Far West, Missouri, to be built up as a city of Zion with a temple but directed the presidency not to go into debt to build the Far West temple as they had when building the Kirtland temple. The issue of JS's and Rigdon's debts was raised again in May when the two petitioned the high council to obtain compensation for their services in the church. Debts continued to loom over JS and Rigdon, and on 8 July 1838 the first payment on a debt totaling over \$4,000 was due to JS's attorneys.
			
			``That day, JS dictated this revelation on tithing---apparently in a church leadership meeting held in Far West. This revelation was the third dictated that day that George W. Robinson copied into JS's journal. A copy of the revelation states that JS dictated it in direct response to the petition, `Lord, show unto thy servents how much thou requirest of the properties of thy people for a Tithing?' The resulting revelation called for the Latter-day Saints to consecrate all of their surplus property and thereafter to pay `one tenth of all their interest annually.' Robinson, who was present at the leadership meeting, may have transcribed the revelation as JS dictated it.
			
			``The revelation was read later that day to a congregation of Latter-day Saints. Over the next few weeks, church members responded to the revelation by consecrating surplus property. According to the 27 July entry in JS's journal, `Some time past the bretheren or saints have come up day after day to consecrate, and to bring their offerings into the store house of the lord, . . . They have come up hither Thus far, according to the ord[e]r of the Dan-Ites.' Officers in the Danite society had attended the leadership meeting in which JS apparently dictated this revelation, and members of the society were now helping gather the consecrated goods. The success in collecting surplus property apparently did not last long. John Corrill, the keeper of the storehouse, recounted in 1839 that the Danites `set out to enforce the law of consecration; but this did not amount to much.' Brigham Young, who was serving in the pro tempore church presidency in Zion, recollected several years later that church members were sparing in what they considered surplus property.
			
			``Robinson made a copy of the revelation in JS's journal, apparently sometime in mid- or late July 1838. The revelation was also copied by other church leaders: Newel K. Whitney made one copy, and Edward Partridge made at least two copies. A comparison of the copies by Robinson, Whitney, and Partridge suggests that one of Partridge's copies most closely represents the wording of the original revelation. This version is featured here. Partridge was present when the revelation was dictated and probably made the featured copy shortly thereafter; the latest possible copying date is 27 May 1840, the day he died.''
			
		% --
		\paragraph{HDDC}
		% --
		
			HDDC 1551--1553 has some extra background here, including somewhat conflicting views (between Joseph Fielding Smith and Hyrum Andrus) about the origin of this revelation.
			
			The more important background is given above in the JSP materials.
					
	% -----
	\subsubsection{Versions}
	% -----
		\label{DC119-V}
		
		See the \href{https://drive.google.com/file/d/1goAvNz8jS4R8slYfMG_db55Ty2z_vmgK/view?usp=sharing}{sections comparison} document.
	
		\begin{compactdesc}
			\item[JSP] \hyperref[EMS-1838-JS-8-JUL-1838-C]{8 July 1838}
			\item[BC] ---
			\item[DC 1835] ---
			\item[DC 1844] Section 107, 430--431
			\item[TS] 5:15 (15 August 1844), 618
			\item[DN] 3:10 (2 April 1853), 37
			\item[MS] 16:12 (25 March 1854), 183
		\end{compactdesc}

	% -----
	\subsubsection{Section 119:1} % *DC119-1 
	% -----
		\label{DC119-1}

		VERILY, thus saith the Lord, I require all their surplus property to be put into the hands of the bishop of my church in Zion,

		% \begin{framed}
		% \scriptsize
			% \begin{compactdesc}
				% \item[JSP] 
				% % \item[BC] 
				% % \item[]
			% \end{compactdesc}
		% \end{framed}

	% -----
	\subsubsection{Section 119:2} % *DC119-2 
	% -----
		\label{DC119-2}

		For the building of mine house, and for the laying of the foundation of Zion and for the priesthood, and for the debts of the Presidency of my Church.

		% \begin{framed}
		% \scriptsize
			% \begin{compactdesc}
				% \item[JSP] 
				% % \item[BC] 
				% % \item[]
			% \end{compactdesc}
		% \end{framed}

	% -----
	\subsubsection{Section 119:3} % *DC119-3 
	% -----
		\label{DC119-3}

		And this shall be the beginning of the tithing of my people.

		% \begin{framed}
		% \scriptsize
			% \begin{compactdesc}
				% \item[JSP] 
				% % \item[BC] 
				% % \item[]
			% \end{compactdesc}
		% \end{framed}

	% -----
	\subsubsection{Section 119:4} % *DC119-4 
	% -----
		\label{DC119-4}

		And after that, those who have thus been tithed shall pay one-tenth of all their \hyperref[INTEREST]{interest}%
			\footnote{%
				% \label{}%
				When you look at the multiple mentions of ``surplus'' in this section, and when you look at the definitions of ``interest'' (\hyperref[INTEREST]{here}), where the most relevant meaning seems to be \#5, ``any surplus advantage,'' it seems like the original tithing commandment did not in fact refer to 10\% of whatever you would call your ``salary'' per year. The emphasis seems rather to be on \textit{surplus} property, whatever that means. This isn't to say that the commandment couldn't be re-interpreted in some other way later, but the salary meaning doesn't seem to be the original one.
				}
		annually; and this shall be a standing law unto them forever, for my holy priesthood, saith the Lord.

		% \begin{framed}
		% \scriptsize
			% \begin{compactdesc}
				% \item[JSP] 
				% % \item[BC] 
				% % \item[]
			% \end{compactdesc}
		% \end{framed}

	% -----
	\subsubsection{Section 119:5} % *DC119-5 
	% -----
		\label{DC119-5}

		Verily I say unto you, it shall come to pass that all those who gather unto the land of Zion shall be tithed of their surplus properties, and shall observe this law, or they shall not be found worthy to abide among you.

		% \begin{framed}
		% \scriptsize
			% \begin{compactdesc}
				% \item[JSP] 
				% % \item[BC] 
				% % \item[]
			% \end{compactdesc}
		% \end{framed}

	% -----
	\subsubsection{Section 119:6} % *DC119-6 
	% -----
		\label{DC119-6}

		And I say unto you, if my people observe not this law, to keep it holy, and by this law sanctify the land of Zion unto me, that my statutes and my judgments may be kept thereon, that it may be most holy, behold, verily I say unto you, it shall not be a land of Zion unto you.%
			\footnote{%
				% \label{}%
				Wherever you are, you make your ``land'' a land of Zion by your actions.
				}

		% \begin{framed}
		% \scriptsize
			% \begin{compactdesc}
				% \item[JSP] 
				% % \item[BC] 
				% % \item[]
			% \end{compactdesc}
		% \end{framed}

	% -----
	\subsubsection{Section 119:7} % *DC119-7 
	% -----
		\label{DC119-7}

		And this shall be an ensample unto all the stakes of Zion.  Even so.  Amen.

		% \begin{framed}
		% \scriptsize
			% \begin{compactdesc}
				% \item[JSP] 
				% % \item[BC] 
				% % \item[]
			% \end{compactdesc}
		% \end{framed}